\documentclass[11pt,a4paper]{article}

\usepackage[T1]{fontenc}
\usepackage[utf8]{inputenc}
\usepackage{lmodern}
\usepackage[margin=1in]{geometry}
\usepackage{amsmath,amssymb,amsthm,mathtools}
\usepackage{booktabs}
\usepackage{array}
\usepackage{tabularx}
\usepackage{float}
\usepackage[hidelinks]{hyperref}
\usepackage{microtype}
\usepackage{enumitem}
\usepackage{siunitx}
\usepackage{xcolor}
\usepackage{tcolorbox}
\tcbuselibrary{skins,breakable}

\sisetup{
    detect-all,
    round-mode = figures,
    round-precision = 10,
    exponent-mode = input,
    exponent-product = \times,
    group-separator = {\,},
    group-minimum-digits = 4
}

\newtheorem{theorem}{Theorem}[section]
\newtheorem{corollary}[theorem]{Corollary}
\newtheorem{lemma}[theorem]{Lemma}
\theoremstyle{definition}
\newtheorem{definition}[theorem]{Definition}
\newtheorem{conjecture}[theorem]{Conjecture}
\newtheorem{proposition}[theorem]{Proposition}
\theoremstyle{remark}
\newtheorem*{remark}{Remark}

\newcommand{\Mpl}{\bar M_{\mathrm{Pl}}}
\newcommand{\cthree}{c_3}
\newcommand{\phitree}{\varphi_{\mathrm{tree}}}
\newcommand{\phiz}{\varphi_0}
\newcommand{\deltatop}{\delta_{\mathrm{top}}}
\newcommand{\bone}{b_1}
\newcommand{\tr}{\operatorname{Tr}}
\newcommand{\rank}{\operatorname{rank}}
\newcommand{\SM}{\mathrm{SM}}
\newcommand{\Hch}{\mathcal H_{\mathrm{ch}}}
\newcommand{\Oadm}{\Omega_{\mathrm{adm}}}
\newcommand{\Cmin}{\mathcal C_{\min}}

\title{%
    \textbf{TFPT v2.9: Admissible Vacua, Minimal Carrier, and Electromagnetic Stationarity}\\[0.5em]
    \Large A Carrier-First Primary Form
}
\author{Stefan Hamann \and Alessandro Rizzo}
\date{\textbf{Draft for Version 2.9 -- 10 March 2026}}

\begin{document}
\maketitle

\begin{abstract}
This document rewrites TFPT v2.8 into a carrier-first primary form. The rigid core is the Minimal Carrier Theorem: admissibility on the orientable double cover together with the seam involution forces the minimal carrier $E_3\oplus E_2$. From this rigid package follow the global Standard Model group $S(U(3)\times U(2))$, one chiral family $S^+=\Lambda^{\mathrm{even}}E$, the carrier norm $\tr_EY^2=5/6$, and the linear and cubic trace cancellations. On the canonical occupancy branch, the admissible chiral multiplicity is
\[
\Oadm = N_{\mathrm{fam}}\dim S^+ = 3\times16 = 48.
\]
The abelian coefficient takes the carrier form
\[
\bone(N_{\mathrm{fam}},N_H)=\frac43 N_{\mathrm{fam}}+\frac1{10}N_H,
\]
so the canonical branch $(N_{\mathrm{fam}},N_H)=(3,1)$ gives $\bone=41/10$. The electromagnetic stationarity equation can then be written entirely in terms of the carrier packet
\[
\Cmin=(E_3\oplus E_2,\;Y,\;S^+,\;\gamma,\;B),
\]
with $\gamma=\tr_EY^2=5/6$ and $B=\rank E_3/\rank E_2=3/2$. Seam descent for $\cthree$, the induced Einstein channel, the E$_8$ reinterpretation, Higgs selection, and discrete phase grammar are retained as structural bridges or dynamical completion targets, but no longer presented as part of the hard carrier core.
\end{abstract}

\begin{tcolorbox}[colback=gray!5, colframe=gray!60, breakable,
    title={\textbf{Status boundary for the v2.9 rewrite}}, fonttitle=\bfseries]
\textbf{1. Hard carrier statements}
\begin{itemize}[leftmargin=1.5em, topsep=3pt, itemsep=2pt]
    \item minimal carrier $E_3\oplus E_2$,
    \item global Standard Model group $S(U(3)\times U(2)) \cong (SU(3)\times SU(2)\times U(1)_Y)/\mathbb Z_6$,
    \item $S^+=\Lambda^{\mathrm{even}}E$ as one chiral family with $\nu^c$,
    \item carrier invariants $\dim S^+=16$, $\gamma=\tr_EY^2=5/6$, and $B=\rank E_3/\rank E_2=3/2$,
    \item linear and cubic trace cancellation,
    \item the general abelian index formula $\bone(N_{\mathrm{fam}},N_H)=\frac43 N_{\mathrm{fam}}+\frac1{10}N_H$.
\end{itemize}

\textbf{2. Structural bridges}
\begin{itemize}[leftmargin=1.5em, topsep=3pt, itemsep=2pt]
    \item three families from the three distinguished cycles $C_1,C_2,C_T$,
    \item $\Oadm=N_{\mathrm{fam}}\dim S^+$ and $\Oadm=48$ on the canonical branch,
    \item $\deltatop = \Oadm\,\cthree^4$ as occupancy transgression,
    \item $\delta_2 = B\gamma\,\Oadm^2\cthree^8$ and therefore $B\gamma=5/4$ on the minimal carrier,
    \item the E$_8$ rereading $\lambda_{E8}=\gamma/\ln(248/60)$.
\end{itemize}

\textbf{3. Dynamical completions}
\begin{itemize}[leftmargin=1.5em, topsep=3pt, itemsep=2pt]
    \item seam descent for $\cthree$,
    \item the relative-$a_2$ / induced-Einstein channel,
    \item the light seam-even Higgs candidate in $E_2$,
    \item the physical realization of the $\mathbb Z_6\rtimes \mathbb Z_2$ phase grammar.
\end{itemize}
\end{tcolorbox}

\tableofcontents

\section{Overview and Claim Map}

The v2.9 rewrite should be read from admissibility to carrier, from carrier to occupancy, and only then from occupancy to closure and phenomenology. The master line is therefore
\[
(\widetilde M,\Gamma,\tau)
\Rightarrow E_3\oplus E_2
\Rightarrow Y,\ S^+,\ \gamma,\ B
\Rightarrow \bone,\ \Oadm
\Rightarrow \text{CFE}
\Rightarrow \alpha.
\]
The main bridge branches are secondary to this line:
\[
\Oadm \Rightarrow \deltatop,\ \delta_2,
\qquad
\gamma \Rightarrow \lambda_{E8},
\qquad
\Delta a_2^{\mathrm{rel}} \Rightarrow \Mpl.
\]
The center of the architecture is therefore not the number $48$ by itself, but the rigid minimal carrier package that later induces $48$ on the canonical occupancy branch.

\begin{table}[H]
\centering
\small
\begin{tabularx}{\linewidth}{@{}>{\raggedright\arraybackslash}p{0.24\linewidth}>{\raggedright\arraybackslash}p{0.14\linewidth}X>{\raggedright\arraybackslash}p{0.12\linewidth}@{}}
\toprule
Claim & Status & Depends on & First derived in section \\
\midrule
Minimal carrier $E_3\oplus E_2$ & hard & admissibility on $\widetilde M$ and seam involution $\tau$ & 3 \\
Global Standard Model group & hard & minimal carrier split & 3 \\
One chiral family $S^+$ & hard & minimal carrier and half-spinor construction & 4 \\
General abelian index $\bone(N_{\mathrm{fam}},N_H)$ & hard & carrier traces plus scalar multiplicity & 4 \\
Canonical occupancy $\Oadm=48$ & structural bridge & family-index bridge and $\dim S^+=16$ & 5 \\
Carrier form of the CFE & structural bridge & $\cthree$, $\Oadm$, $\bone$, $B$, $\gamma$ & 6 \\
Induced Einstein channel & dynamical target & relative spectral completion and canonical occupancy & 8 \\
E$_8$ rereading & downstream reinterpretation & carrier norm $\gamma$ & 9 \\
\bottomrule
\end{tabularx}
\caption{Claim map for the carrier-first v2.9 presentation. Appendix~A records a compact v2.8$\to$v2.9 dictionary.}
\end{table}

\section{Admissibility and Seam Data}

The ambient v2.9 datum may be packaged as
\[
\mathfrak T_{2.9} = \bigl(\widetilde M \to M,\ \Gamma,\ \tau,\ \chi\bigr),
\]
but the hard carrier core uses only $(\widetilde M,\Gamma,\tau)$. The field $\chi$ belongs to the absolute-scale and gravity channels discussed later in Section~8.

\subsection{Seam data recalled from v2.8}

The seam information needed by the carrier rewrite is already present in v2.8. In the minimal gluing class one has
\[
U_\Gamma(\theta)=e^{i\theta},
\qquad
\mathrm{SF}(U_\Gamma)=1,
\qquad
\Delta_\Gamma=2\pi.
\]
On the orientable double cover the distinguished cycles are
\[
C_1,\qquad C_2,\qquad C_T,
\]
where $C_T$ is the seam cycle in the glued description and corresponds to $\Gamma_+\cup\Gamma_-$ in the cut description. These data should appear near the front of the paper because Sections~3--6 all depend on them.

\begin{tcolorbox}[colback=blue!3, colframe=blue!50!black,
    title={\textbf{What $\chi$ does and does not do}}, fonttitle=\bfseries\small]
$\chi$ belongs to the absolute-scale and gravitational channels. It does \emph{not} enter the hard derivation of the minimal carrier, the global Standard Model group, the one-family half-spinor, or the carrier norm $\tr_EY^2=5/6$. Only after the carrier packet has been fixed does $\chi$ enter the relative spectral stationarity problem and the induced-gravity target.
\end{tcolorbox}

\section{The Minimal Carrier Theorem}

\subsection{Uniqueness of the \texorpdfstring{$3+2$}{3+2} split}

\begin{theorem}[Minimal carrier theorem]
Let $E=E_b\oplus E_s$ be a split carrier with traceless diagonal generator
\[
Y_{b,s}=\mathrm{diag}\!\left(-\frac1b,\ldots,-\frac1b,\frac1s,\ldots,\frac1s\right).
\]
If one simultaneously requires
\[
\dim \Lambda^{\mathrm{even}}E = 16
\qquad\text{and}\qquad
\tr_E Y_{b,s}^2 = \frac56,
\]
then necessarily $\{b,s\}=\{3,2\}$.
\end{theorem}

\begin{proof}
For a rank-$g$ carrier with $g=b+s$, the positive half-spinor has dimension
\[
\dim \Lambda^{\mathrm{even}}E = 2^{g-1}.
\]
Demanding one chiral family gives $2^{g-1}=16$, hence $g=5$ and therefore $b+s=5$.

For the traceless generator above one has
\[
\tr_E Y_{b,s}^2 = b\left(\frac{1}{b^2}\right)+s\left(\frac{1}{s^2}\right)
= \frac1b+\frac1s
= \frac{b+s}{bs}
= \frac{5}{bs}.
\]
Imposing $\tr_EY_{b,s}^2=5/6$ yields $bs=6$. Together with $b+s=5$, this forces $(b,s)=(3,2)$ up to order.
\end{proof}

Hence from now on we set
\[
E=E_3\oplus E_2,
\qquad
P_3:=\frac{1-\tau}{2},
\qquad
P_2:=\frac{1+\tau}{2}.
\]

\begin{proposition}[Hypercharge as an affine form of the seam involution]
For the minimal carrier,
\[
Y=-\frac13 P_3+\frac12 P_2
=\mathrm{diag}\!\left(-\frac13,-\frac13,-\frac13,\frac12,\frac12\right),
\]
and equivalently
\[
Y=\frac{1}{12}\mathbf 1+\frac{5}{12}\tau.
\]
\end{proposition}

\begin{proof}
Using $P_3=(1-\tau)/2$ and $P_2=(1+\tau)/2$ gives
\[
-\frac13 P_3+\frac12 P_2
= -\frac16(1-\tau)+\frac14(1+\tau)
= \frac{1}{12}\mathbf 1+\frac{5}{12}\tau.
\]
\end{proof}

\begin{theorem}[Global Standard Model theorem]
The determinant-preserving connected stabilizer of the minimal carrier split is
\[
\mathrm{Aut}_0(E_3\oplus E_2,\det=1)
= S(U(3)\times U(2))
\cong \frac{SU(3)\times SU(2)\times U(1)_Y}{\mathbb Z_6}.
\]
\end{theorem}

\begin{proof}
Automorphisms preserving the split and the determinant condition act independently on $E_3$ and $E_2$ with unit determinant on the combined carrier. This gives $S(U(3)\times U(2))$. The covering map from $SU(3)\times SU(2)\times U(1)_Y$ to $S(U(3)\times U(2))$ has the canonical $\mathbb Z_6$ subgroup as kernel, yielding the stated global Standard Model form.
\end{proof}

Set
\[
S^+ := \Lambda^{\mathrm{even}}(E_3\oplus E_2).
\]

\begin{proposition}[Carrier invariants]
For the minimal carrier,
\[
\gamma := \tr_EY^2 = \frac56,
\qquad
B := \frac{\rank E_3}{\rank E_2} = \frac32,
\qquad
\dim S^+ = 16.
\]
\end{proposition}

\begin{proof}
The first identity is
\[
\tr_EY^2 = 3\left(\frac19\right)+2\left(\frac14\right)=\frac13+\frac12=\frac56.
\]
The second is immediate from the split ranks. The third follows from
\[
\dim S^+ = \dim \Lambda^{\mathrm{even}}E = 2^{5-1}=16.
\]
\end{proof}

\begin{definition}[Minimal carrier packet]
The rigid carrier packet of v2.9 is
\[
\Cmin := (E_3\oplus E_2,\;Y,\;S^+,\;\gamma,\;B).
\]
\end{definition}

\section{Carrier Matter, Trace Identities, and the Abelian Index}

\subsection{One chiral family from the positive half-spinor}

\begin{theorem}[Carrier matter theorem]
For the minimal carrier $E=E_3\oplus E_2$, the positive half-spinor
\[
S^+ := \Lambda^{\mathrm{even}}(E_3\oplus E_2)
\]
decomposes as
\[
(1,1)_0
\oplus
(3,2)_{1/6}
\oplus
(\bar 3,1)_{-2/3}
\oplus
(\bar 3,1)_{1/3}
\oplus
(1,2)_{-1/2}
\oplus
(1,1)_1.
\]
Thus $S^+$ carries exactly one Standard Model family with a right-handed neutrino:
\[
\nu^c,\ Q,\ u^c,\ d^c,\ L,\ e^c.
\]
\end{theorem}

\begin{proof}
The even exterior powers split as
\[
\Lambda^{\mathrm{even}}E
=
\Lambda^0 E \oplus \Lambda^2 E \oplus \Lambda^4 E.
\]
Now
\[
\Lambda^0 E = (1,1)_0,
\]
\[
\Lambda^2 E = \Lambda^2 E_3 \oplus (E_3\otimes E_2)\oplus \Lambda^2 E_2
= (\bar 3,1)_{-2/3}\oplus (3,2)_{1/6}\oplus (1,1)_1,
\]
and $\Lambda^4E \cong E^\ast$ because $\det E$ is trivial in the $S(U(3)\times U(2))$ form, so
\[
\Lambda^4E \cong (\bar 3,1)_{1/3}\oplus (1,2)_{-1/2}.
\]
Collecting the terms gives the stated multiplets.
\end{proof}

\subsection{Trace identities}

\begin{theorem}[General half-spinor quadratic trace identity]
Let $Y$ be a traceless diagonal generator on a rank-$g$ carrier $E$. Then on the positive half-spinor
\[
S^+ = \Lambda^{\mathrm{even}}E
\]
one has
\[
\tr_{S^+}Y^2 = 2^{g-3}\,\tr_EY^2.
\]
\end{theorem}

\begin{proof}
Write the diagonal weights of $Y$ as $q_i$ with $\sum_i q_i=0$. Then
\[
\tr_{S^+}Y^2
=
\sum_{S\ \mathrm{even}}
\left(\sum_{i\in S} q_i\right)^2.
\]
Each $q_i^2$ appears in $2^{g-2}$ even subsets and each pair $q_iq_j$ appears in $2^{g-3}$ even subsets. Using
\[
\left(\sum_i q_i\right)^2=0
\]
shows that the cross terms cancel half the self-term contribution, leaving
\[
\tr_{S^+}Y^2
=
2^{g-3}\sum_i q_i^2
=
2^{g-3}\tr_EY^2.
\]
\end{proof}

\begin{corollary}[Minimal carrier quadratic trace]
For the minimal carrier,
\[
\tr_{S^+}Y^2 = 2^{5-3}\tr_EY^2 = 4\cdot\frac56 = \frac{10}{3}.
\]
\end{corollary}

\begin{theorem}[Linear and cubic trace cancellation]
On the one-family carrier spectrum $S^+$ one has
\[
\tr_{S^+}Y = 0,
\qquad
\tr_{S^+}Y^3 = 0.
\]
\end{theorem}

\begin{proof}
Using the explicit family decomposition above, the hypercharges are
\[
0,\quad \frac16,\quad -\frac23,\quad \frac13,\quad -\frac12,\quad 1,
\]
with multiplicities $1,6,3,3,2,1$. Hence
\[
\tr_{S^+}Y
= 6\cdot\frac16+3\cdot\!\left(-\frac23\right)+3\cdot\frac13+2\cdot\!\left(-\frac12\right)+1=0,
\]
and
\[
\tr_{S^+}Y^3
= 6\cdot\left(\frac16\right)^3
+3\cdot\left(-\frac23\right)^3
+3\cdot\left(\frac13\right)^3
+2\cdot\left(-\frac12\right)^3
+1^3=0.
\]
\end{proof}

\subsection{General abelian index}

\begin{theorem}[General carrier formula for the abelian coefficient]
For a branch with $N_{\mathrm{fam}}$ copies of the carrier family and $N_H$ light weak doublets from $E_2$, one has
\[
b_Y(N_{\mathrm{fam}},N_H)
:= \frac23 N_{\mathrm{fam}}\tr_{S^+}Y^2 + \frac13 N_H\tr_{E_2}Y^2.
\]
For the minimal carrier this becomes
\[
b_Y(N_{\mathrm{fam}},N_H)
= \frac23 N_{\mathrm{fam}}\cdot\frac{10}{3} + \frac13 N_H\cdot\frac12
= \frac{20}{9}N_{\mathrm{fam}}+\frac16 N_H.
\]
Therefore
\[
\bone(N_{\mathrm{fam}},N_H)
= \frac35\,b_Y(N_{\mathrm{fam}},N_H)
= \frac43 N_{\mathrm{fam}}+\frac1{10}N_H.
\]
\end{theorem}

\begin{corollary}[Canonical carrier branch]
For the canonical branch $(N_{\mathrm{fam}},N_H)=(3,1)$,
\[
b_Y=\frac{41}{6},
\qquad
\bone=\frac{41}{10}.
\]
\end{corollary}

\begin{remark}
The symbol $\nu^c$ is already part of the hard packet $S^+$. Removing it does not define a different carrier; it is an external truncation of the family packet and is treated as such only in Section~7.
\end{remark}

\section{Boundary Multiplicity and Admissible Occupancy}

\subsection{Family-index bridge}

\begin{conjecture}[Family-index bridge]
Let each distinguished cycle $C_i$ carry an integer family index $I_{\mathrm{fam}}(C_i)$. Then the v2.8 block coefficient may be organized as
\[
k_B = \frac12\sum_{i=1}^{3} I_{\mathrm{fam}}(C_i)\,I_1(B).
\]
\end{conjecture}

Combined with the v2.8 block formula
\[
k_B=\frac32\,I_1(B),
\]
this gives
\[
\sum_{i=1}^{3} I_{\mathrm{fam}}(C_i)=3.
\]
The equal-one assignment
\[
I_{\mathrm{fam}}(C_1)=I_{\mathrm{fam}}(C_2)=I_{\mathrm{fam}}(C_T)=1
\]
is the minimal symmetric branch and yields
\[
N_{\mathrm{fam}}=3.
\]

\subsection{Admissible occupancy and topological surcharge}

\begin{definition}[Admissible occupancy]
The admissible chiral occupancy is
\[
\Oadm := \tr_{\Hch}\mathbf 1 = N_{\mathrm{fam}}\dim S^+.
\]
\end{definition}

On the canonical branch one therefore has
\[
\Oadm = 3\times 16 = 48.
\]

\begin{conjecture}[Occupancy transgression]
The topological surcharge should be written as
\[
\deltatop = \Oadm\,\cthree^4.
\]
\end{conjecture}

For the canonical branch this becomes
\[
\deltatop = 48\cthree^4 = \frac{3}{256\pi^4} = 3\cdot 16\,\cthree^4.
\]
In this reading the factor $3$ comes from the weighted boundary geometry while the factor $16$ comes from the minimal carrier packet itself.

\subsection{Seam descent normalization of \texorpdfstring{$\cthree$}{c3}}

The normalization of $\cthree$ is best presented in two stages.

\textbf{Stage 1: hard seam geometry.}
On the spin lift of the minimal seam loop one has
\[
\widetilde U_\Gamma(2\pi)=-1,
\qquad
\widetilde U_\Gamma(4\pi)=1.
\]
So the spin-lifted loop closes only after $4\pi$.

\textbf{Stage 2: conditional axionic identification.}
If the admissible carrier modulus is identified with the axionic angle variable, write
\[
a=\frac{\theta}{4\pi}.
\]
Then
\[
\mathcal L_\theta = \frac{\theta}{32\pi^2}F\widetilde F
= \frac{1}{8\pi}aF\widetilde F,
\]
hence
\[
\cthree = \frac{1}{8\pi}.
\]
The geometry and the final normalization are thus explicitly separated: the $4\pi$ closure is geometric, while the identification of the carrier modulus with the axionic angle remains the extra conditional step.

\section{Electromagnetic Stationarity and the Carrier CFE}

\subsection{Carrier form of the stationary equation}

The electromagnetic closure channel should first be written on a general admissible branch. Define
\[
\phiz(\alpha)
= \frac{1}{6\pi}
+ \Oadm\,\cthree^4 e^{-2\alpha}
+ B\gamma\,\Oadm^2\cthree^8 e^{-4\alpha}
+ \cdots.
\]
Then the carrier form of the CFE is
\[
\alpha^3
-2\cthree^3\alpha^2
-8\cthree^6\,\bone(N_{\mathrm{fam}},N_H)
\ln\!\bigl(\phiz(\alpha)^{-1}\bigr)
=0.
\]
This is the preferred v2.9 presentation because the coefficients are now grouped by carrier data rather than as parallel numerical motifs.

\begin{proposition}[Canonical carrier branch]
For
\[
(N_{\mathrm{fam}},N_H,\nu^c)=(3,1,\text{present}),
\]
one has
\[
\Oadm=48,
\qquad
\bone=\frac{41}{10},
\qquad
B\gamma=\frac32\cdot\frac56=\frac54.
\]
Hence
\[
\phiz(\alpha)
= \frac{1}{6\pi}
+ 48\cthree^4 e^{-2\alpha}
+ \frac54\,48^2\cthree^8 e^{-4\alpha}
+ \cdots.
\]
\end{proposition}

\begin{remark}
The identity
\[
\frac54 = B\gamma
\]
is the compressed carrier reading of the v2.8 two-defect weight.
\end{remark}

Using
\[
\cthree=\frac{1}{8\pi},
\qquad
\phitree=\frac{1}{6\pi},
\qquad
\deltatop=48\cthree^4,
\]
the positive stationary root remains
\[
\alpha^{-1}(0)=\num{137.03599921615843}.
\]
The point of v2.9 is therefore not to move the number, but to reduce the number of primitive inputs behind it.

\begin{table}[H]
\centering
\small
\begin{tabularx}{\linewidth}{@{}lXl@{}}
\toprule
Quantity & Carrier meaning & Value \\
\midrule
$\cthree$ & canonical seam-descent normalization & $\frac{1}{8\pi}$ \\
$\Oadm$ & admissible chiral occupancy & $48$ \\
$\gamma$ & carrier norm $\tr_EY^2$ & $\frac56$ \\
$B$ & carrier rank ratio $\rank E_3/\rank E_2$ & $\frac32$ \\
$B\gamma$ & compressed two-defect weight & $\frac54$ \\
$\bone$ & canonical abelian coefficient & $\frac{41}{10}$ \\
$\alpha^{-1}(0)$ & positive root of the carrier CFE & \num{137.03599921615843} \\
\bottomrule
\end{tabularx}
\caption{Numerical summary of the canonical carrier branch.}
\end{table}

\subsection{Algebraic corollaries of the carrier kernel}

Evaluated on the stationary root of the canonical branch, several previously separate formulas become algebraic corollaries of the same carrier package:
\[
g_{a\gamma\gamma}=-4\cthree=-\frac{1}{2\pi},
\qquad
\beta_{\mathrm{rad}}=\frac{\phiz}{4\pi},
\]
\[
\lambda_C=\sqrt{\phiz(1-\phiz)},
\qquad
\sin^2\theta_{13}=\phiz e^{-\gamma}.
\]
The structural gain is not that these quantities are numerically new, but that they are now visibly downstream of $\Cmin$ and the occupancy branch.

\subsection{Conditional \texorpdfstring{$\chi$}{chi}-channel corollaries}

If the $\chi$ channel and the induced Einstein channel are successfully completed, then the inflationary block can be written as
\[
\frac{M}{\chi_0}=\frac{2}{\sqrt{\pi}}\,\cthree^4,
\qquad
\frac{M}{\Mpl}=\sqrt{8\pi}\,\cthree^4,
\]
and therefore
\[
A_s=\frac{N^2}{24\pi^2}\left(\frac{M}{\Mpl}\right)^2
= \frac{N^2}{3\pi}\cthree^8.
\]
For $N=56$ this yields
\[
A_s=\num{2.09019136432946e-9}.
\]
These results are kept separate from the hard carrier corollaries because they depend on the success of the relative-$\chi$ and gravity channels.

\section{Discrete Exclusion Structure}

The carrier-centered rewrite is strongest when it excludes wrong branches, not when it celebrates an isolated number. The relevant deformations fall into three logical classes.

\begin{table}[H]
\centering
\small
\begin{tabularx}{\linewidth}{@{}l l X l@{}}
\toprule
Logical type & Deformation & Carrier/CFE consequence & Result \\
\midrule
Boundary multiplicity & $N_{\mathrm{fam}}=2$ & $\Oadm=32$ and $\bone=\frac{83}{30}$ in the carrier CFE & $\alpha^{-1}(0)=\num{156.094045457}$ \\
Boundary multiplicity & $N_{\mathrm{fam}}=4$ & $\Oadm=64$ and $\bone=\frac{163}{30}$ in the carrier CFE & $\alpha^{-1}(0)=\num{124.835054800}$ \\
Scalar content & $N_H=2$ & $\bone=\frac{21}{5}$ with the same carrier family packet & $\alpha^{-1}(0)=\num{135.945991742}$ \\
External truncation & no $\nu^c$ & the family packet is externally reduced from $16$ to $15$ states, so $\Oadm=45$ and the hard carrier theorem no longer applies unchanged & $\alpha^{-1}(0)=\num{137.033841585}$ \\
\bottomrule
\end{tabularx}
\caption{Discrete exclusion structure in the carrier reading. The ``no $\nu^c$'' row is an external truncation of the hard family packet, not a normal carrier alternative.}
\end{table}

In addition, a light seam-even colored triplet zero mode would require a light seam-even bosonic mode in $E_3$. The minimal carrier picture does not provide such a mode, so the colored partner is not on equal footing with the Higgs candidate in $E_2$.

\section{Relative Spectral Completion Targets}

\begin{tcolorbox}[colback=gray!5, colframe=gray!60, breakable,
    title={\textbf{Status: dynamical target, not part of the hard carrier core}}, fonttitle=\bfseries]
The relative spectral action is the natural place to unify the $\alpha$-, $\chi$-, and metric channels, but the heat-kernel completion is not yet theorem-hard. In v2.9 it should therefore appear as a completion target rather than as part of the rigid carrier packet.
\end{tcolorbox}

The natural completion ansatz is
\[
S_{\mathrm{rel}}
= \tr f(D/\chi)
- \tr f(D_{\mathrm{ref}}/\chi)
+ \frac{i\pi}{2}\Delta\eta_\Gamma
+ S_{\mathrm{defect}}.
\]
The three stationarity channels are then
\[
\frac{\delta S_{\mathrm{rel}}}{\delta\alpha}=0,
\qquad
\frac{\delta S_{\mathrm{rel}}}{\delta\chi}=0,
\qquad
\frac{\delta S_{\mathrm{rel}}}{\delta e}=0,
\]
with an independent $\omega$ variation left to a later Palatini/Einstein--Cartan completion.

\subsection{Why the denominator 192 pi squared is plausible}

For a Dirac-type Laplace operator written as
\[
D^2=-\nabla^2+\mathcal E,
\]
the local four-dimensional heat-kernel coefficient is
\[
a_2(D^2)=\frac{1}{16\pi^2}\int d^4x\sqrt{-g}\,
\tr\!\left(\mathcal E+\frac16 R\right).
\]
For a Dirac-type curvature coupling one has $\mathcal E\sim R/4+\cdots$, so the curvature part scales like
\[
\frac{1}{16\pi^2}\left(\frac14+\frac16\right)
\int d^4x\sqrt{-g}\,\tr(\mathbf 1)\,R
=
\frac{5}{192\pi^2}
\int d^4x\sqrt{-g}\,\tr(\mathbf 1)\,R.
\]
After relative admissible-minus-reference subtraction and TFPT normalization, an Einstein-channel coefficient with denominator $192\pi^2$ is therefore natural. What remains open is the exact numerator and the precise sector trace selected by the full relative operator.

\begin{conjecture}[Induced Einstein channel]
The relative spectral completion should organize the metric channel as
\[
S_{\mathrm{EH}}^{\mathrm{rel}}
=
\frac{\chi^2}{192\pi^2}\,\tr_{\Hch}\mathbf 1
\int d^4x\sqrt{-g}\,R.
\]
On the canonical occupancy branch, where $\tr_{\Hch}\mathbf 1=\Oadm=48$, this gives
\[
S_{\mathrm{EH}}^{2.9}
= \int d^4x\sqrt{-g}\,\frac{\chi^2}{4\pi^2}R
= \int d^4x\sqrt{-g}\,\frac{\Mpl^2}{2}R,
\]
so that
\[
\Mpl^2=\frac{\chi_0^2}{2\pi^2}.
\]
\end{conjecture}

Using the observed reduced Planck mass only as a calibration check gives
\[
\chi_0 = \sqrt2\,\pi\,\Mpl \approx \num{1.08184e19}\ \mathrm{GeV}.
\]
The v2.9 claim is not that this is already hard, but that the relative-$a_2$ channel is the correct place to make it hard.

\section{Downstream Reinterpretations: \texorpdfstring{E$_8$}{E8}, Higgs, Flavor}

\subsection{\texorpdfstring{E$_8$}{E8} as a downstream carrier rereading}

The strongest E$_8$ rereading to retain is
\[
\lambda_{E8}=\frac{\gamma}{\ln(248/60)},
\qquad
\gamma=\tr_EY^2=\frac56.
\]
The E$_8$ sector is downstream of the carrier norm, not part of the hard minimal carrier core. One suggestive but non-rigid echo is $60=48+12=\Oadm+\dim\mathfrak g_{\SM}$.\footnote{This identity is numerically suggestive, but not theorem-level and should remain secondary or appendix-level.}

\subsection{Higgs as a conditional seam-even mode in \texorpdfstring{$E_2$}{E2}}

If a light seam-even bosonic zero mode exists in $E_2$, then the minimal carrier provides a natural Higgs candidate
\[
\tau H = +H,
\qquad
H\in \Gamma(E_2)\cong (1,2)_{1/2}.
\]
This is the correct representation-theoretic slot for a Standard Model Higgs doublet. The existence of the light mode itself remains part of the dynamical completion and should not be overstated as a theorem of the hard carrier core.

\subsection{Flavor and the phase grammar}

The hard group-theoretic statement is
\[
G_{\mathrm{SM}}^{\mathrm{glob}}
=
S(U(3)\times U(2))
\cong \frac{SU(3)\times SU(2)\times U(1)_Y}{\mathbb Z_6}.
\]
The separate, still dynamical question is whether the observed CKM, PMNS, and Majorana phases actually realize the associated discrete grammar.

\begin{conjecture}[Phase grammar]
The global quotient data suggest a distinguished six-point phase lattice
\[
\delta \in \left\{
0,\frac{\pi}{3},\frac{2\pi}{3},\pi,\frac{4\pi}{3},\frac{5\pi}{3}
\right\}\mod 2\pi.
\]
With CP acting by inversion,
\[
\delta \mapsto -\delta \mod 2\pi,
\]
the resulting phase grammar is naturally
\[
\mathbb Z_6\rtimes \mathbb Z_2.
\]
\end{conjecture}

In v2.9 the quotient theorem should therefore be stated hard, while the physical realization of the phase grammar should remain explicitly conjectural.

\section{Conclusion}

The rigid core of v2.9 is not a number but a theorem:
\[
\boxed{
\text{admissibility on the orientable double cover and the seam involution}
}
\]
\[
\boxed{
\text{force the minimal carrier }E_3\oplus E_2.
}
\]
From this follow the global Standard Model group, one chiral family, the carrier norm $\tr_EY^2=5/6$, and the general abelian index formula. Only on the canonical occupancy branch does the same rigid package produce
\[
\Oadm = N_{\mathrm{fam}}\dim S^+ = 3\times 16 = 48.
\]
The center is therefore the minimal carrier. The number $48$ is the first large trace of that carrier packet in an admissible vacuum, not the foundation of the architecture.

\appendix

\section{v2.8 to v2.9 Dictionary}

\begin{table}[H]
\centering
\small
\begin{tabularx}{\linewidth}{@{}lX@{}}
\toprule
v2.8 language & v2.9 carrier reading \\
\midrule
$g=5$ & $\gamma=\frac56$ together with $B=\frac32$ on the minimal carrier packet \\
$48$ & $\Oadm=N_{\mathrm{fam}}\dim S^+$ on the canonical occupancy branch \\
$\bone=\frac{41}{10}$ & $\bone(N_{\mathrm{fam}},N_H)=\frac43 N_{\mathrm{fam}}+\frac1{10}N_H$ evaluated at $(3,1)$ \\
$\delta_2=\frac54\,\deltatop^2$ & $\delta_2 = B\gamma\,\Oadm^2\cthree^8$ with $B\gamma=\frac54$ \\
$\gamma(0)=\frac56$ in E$_8$ language & $\gamma=\tr_EY^2$ as a carrier invariant \\
\bottomrule
\end{tabularx}
\caption{Compact dictionary from the v2.8 closure language to the v2.9 carrier language.}
\end{table}

\section{Family-Index Bridge}

The v2.8 block relation
\[
k_B=\frac32\,I_1(B)
\]
already shows that family multiplicity is tied to the three distinguished cycles of the orientable double cover. Writing
\[
k_B=\frac12\sum_{i=1}^{3} I_{\mathrm{fam}}(C_i)\,I_1(B)
\]
makes the same factorization explicit and immediately yields
\[
\sum_{i=1}^{3} I_{\mathrm{fam}}(C_i)=3.
\]
The equal-one assignment is the minimal symmetric branch and is the natural source of $N_{\mathrm{fam}}=3$.

\section{Occupancy Transgression}

The canonical bridge is
\[
\Oadm = N_{\mathrm{fam}}\dim S^+ = 3\cdot 16 = 48.
\]
With $\cthree=1/(8\pi)$ this gives
\[
\Oadm\,\cthree^4
= 48\left(\frac{1}{8\pi}\right)^4
= \frac{3}{256\pi^4}
= \deltatop.
\]
This is the precise sense in which the topological surcharge is re-read as admissible occupancy transgression.

\section{Seam Descent Normalization}

The seam descent logic is intentionally separated into two pieces:
\begin{enumerate}[leftmargin=1.5em, topsep=3pt, itemsep=2pt]
    \item the spin-lifted seam loop closes only after $4\pi$, so $\widetilde U_\Gamma(2\pi)=-1$ and $\widetilde U_\Gamma(4\pi)=1$;
    \item if the admissible carrier modulus is identified with the axionic angle variable, then $a=\theta/(4\pi)$ and therefore $\cthree=1/(8\pi)$.
\end{enumerate}
The first step is geometric. The second is the explicit carrier-to-axion identification.

\section{Relative \texorpdfstring{$a_2$}{a2} Sketch}

The relative Einstein target rests on the standard four-dimensional heat-kernel structure
\[
a_2(D^2)=\frac{1}{16\pi^2}\int d^4x\sqrt{-g}\,
\tr\!\left(\mathcal E+\frac16 R\right).
\]
For Dirac-type curvature coupling $\mathcal E\sim R/4+\cdots$, so curvature terms naturally carry a denominator of the form $192\pi^2$. The open step is not whether such a denominator is plausible, but which relative subtraction and sector trace produce the final canonical coefficient
\[
\frac{\chi^2}{192\pi^2}\,\tr_{\Hch}\mathbf 1.
\]

\section{Numerical Branch Table}

\begin{table}[H]
\centering
\small
\begin{tabularx}{\linewidth}{@{}>{\raggedright\arraybackslash}Xcccc@{}}
\toprule
Branch & $N_{\mathrm{fam}}$ & $N_H$ & $\Oadm$ & $\alpha^{-1}(0)$ \\
\midrule
canonical carrier branch & $3$ & $1$ & $48$ & \num{137.03599921615843} \\
boundary multiplicity deformation & $2$ & $1$ & $32$ & \num{156.094045457} \\
boundary multiplicity deformation & $4$ & $1$ & $64$ & \num{124.835054800} \\
scalar-content deformation & $3$ & $2$ & $48$ & \num{135.945991742} \\
external truncation (no $\nu^c$) & $3$ & $1$ & $45$ & \num{137.033841585} \\
\bottomrule
\end{tabularx}
\caption{Representative branches in the carrier CFE. The ``no $\nu^c$'' row is not a normal carrier branch but an external truncation of the hard family packet.}
\end{table}

\end{document}
